\section{Tổng quan kiến trúc}

Hệ thống OISM áp dụng Kiến trúc Sạch nhằm tách biệt quy tắc nghiệp vụ khỏi hạ tầng và giao diện.
Cách tổ chức này giúp hệ thống dễ kiểm thử, thay đổi công nghệ hạ tầng và bảo trì lâu dài.

Kiến trúc được chia thành các lớp chính:

\begin{itemize}
    \item \textbf{Miền nghiệp vụ}: thực thể, quy tắc và logic cốt lõi như đơn hàng, tồn kho, giữ hàng và giá vốn;
    \item \textbf{Ứng dụng}: các ca sử dụng, dịch vụ ứng dụng, giao dịch và điều phối nghiệp vụ;
    \item \textbf{Hạ tầng}: Entity Framework Core, cơ sở dữ liệu, xác thực, tác vụ nền và các tích hợp bên ngoài;
    \item \textbf{Trình bày}: Web API, giao diện quản trị và ứng dụng POS PWA.
\end{itemize}

\section{Các thành phần logic}

Các thành phần chính của hệ thống gồm quản lý danh tính, khách hàng và chi nhánh, danh mục
sản phẩm, tồn kho, tính giá vốn, trung tâm đơn hàng, giữ hàng chống bán vượt, POS, báo cáo,
thông báo thời gian thực và xử lý tác vụ nền.

Web API đóng vai trò điểm truy cập trung tâm. Các yêu cầu nghiệp vụ được chuyển đến lớp ứng
dụng; lớp ứng dụng gọi các dịch vụ miền và kho dữ liệu thông qua các giao diện được định nghĩa
ở lớp bên trong.

\section{Thiết kế SaaS đa khách hàng}

Hệ thống sử dụng mô hình cơ sở dữ liệu dùng chung, trong đó dữ liệu thuộc khách hàng có trường
định danh \texttt{TenantId}. Mọi truy vấn và thao tác ghi đối với dữ liệu thuộc khách hàng phải
được kiểm soát theo ngữ cảnh khách hàng hiện tại.

Entity Framework Core có thể sử dụng Bộ lọc truy vấn toàn cục để giảm nguy cơ truy cập nhầm
dữ liệu giữa các khách hàng. Các cột \texttt{TenantId}, mã SKU, chi nhánh và các trường thường
được tìm kiếm phải có chỉ mục phù hợp.

\section{Công nghệ dự kiến}

\begin{itemize}
    \item .NET 8 Web API cho dịch vụ phía máy chủ;
    \item Entity Framework Core cho truy cập dữ liệu;
    \item cơ sở dữ liệu quan hệ cho dữ liệu nghiệp vụ;
    \item PWA cho giao diện POS;
    \item SignalR cho cập nhật thời gian thực;
    \item Hangfire cho công việc nền;
    \item Docker và Docker Compose cho đóng gói và triển khai;
    \item GitHub Actions cho tích hợp và triển khai liên tục.
\end{itemize}

\section{Kiến trúc giao diện}

Giao diện quản trị phục vụ cấu hình, danh mục, tồn kho, đơn hàng và báo cáo. Giao diện POS
được tối ưu cho thao tác nhanh, tìm kiếm SKU và thanh toán tại quầy. Hai giao diện sử dụng
cùng Web API và các quy tắc nghiệp vụ phía máy chủ.

\section{Nhất quán tồn kho}

Sổ cái tồn kho chỉ ghi thêm là nguồn dữ liệu có khả năng truy vết đối với các biến động. Bên
cạnh đó, hệ thống có thể duy trì bảng số dư tồn kho để đọc nhanh. Mọi thao tác có khả năng làm
thay đổi số dư cuối cùng phải được bảo vệ bằng giao dịch cơ sở dữ liệu và cơ chế khóa thích hợp.

Khi nhiều đơn hàng cùng yêu cầu một SKU, giao dịch phải kiểm tra số lượng có thể bán trên dữ
liệu được khóa trước khi tăng số lượng giữ. Cách làm này ngăn hai giao dịch cùng nhìn thấy một
số lượng tồn giống nhau rồi cùng xác nhận vượt quá hàng thực tế.

\section{Kiến trúc bảo mật}

Xác thực sử dụng mã truy cập JWT và mã làm mới. Mật khẩu phải được lưu dưới dạng giá trị băm
bằng thuật toán phù hợp. Phân quyền được thực hiện ở cả mức vai trò và phạm vi khách hàng.

Thông tin bí mật và cấu hình theo môi trường phải được tách khỏi mã nguồn, không đưa thông tin
xác thực vào kho Git. Giao tiếp giữa các thành phần cần được bảo vệ bằng HTTPS.

\section{Kiến trúc triển khai}

Hệ thống được đóng gói thành các vùng chứa để API, cơ sở dữ liệu, giao diện và các thành phần
hỗ trợ có thể triển khai nhất quán. Docker Compose được dùng cho môi trường trình diễn; quy
trình CI/CD chịu trách nhiệm xây dựng, kiểm thử và đóng gói phiên bản.

\section{Mục tiêu chất lượng}

Kiến trúc được đánh giá theo các tiêu chí:

\begin{itemize}
    \item tính nhất quán và khả năng truy vết của tồn kho;
    \item khả năng chống bán vượt khi có giao dịch đồng thời;
    \item khả năng cô lập dữ liệu đa khách hàng;
    \item hiệu năng tìm kiếm và ghi đơn hàng;
    \item khả năng kiểm thử và bảo trì;
    \item khả năng triển khai có thể tái lập.
\end{itemize}
